\subsection{Orientaciones de soporte completo y equivariancia reticular}
\label{paper:24-orientaciones}

El enumerador de pesos de \eqref{exc:enumerador} contiene
\(24y^{12}\). Por tanto, el conjunto
\[
 C_W^\times:=
 \{u\in C_W:u_b\in\{1,2\}\text{ para todo }b\}
\]
posee 24 elementos. Para cada \(u\in C_W^\times\), definimos
\[
 P_b(u)=
 \begin{cases}
  \mathsf R,&u_b=1,\\
  I,&u_b=2,
 \end{cases}
 \qquad
 g_u:=\bigoplus_{b=0}^{11}P_b(u)\mathsf C P_b(u)^{-1}.
\]

\begin{proposition}[Robustez de las orientaciones de soporte total]
\label{paper:robustez-orientaciones}
Para toda \(u\in C_W^\times\),
\[
 \frac{g_uv-v}{3},
 \quad
 \frac{g_u^{-1}v-v}{3}
 \in N_0.
\]
En particular, cada \(g_u\) preserva el vecino construido con la
orientación transportada.
\end{proposition}

\begin{proof}
En coordenadas de raíces simples,
\[
 \mathsf C\rho-\rho=(-2,-1),
 \qquad
 \mathsf C^{-1}\rho-\rho=(-1,-2).
\]
Las clases discriminantes de los dos desplazamientos son \(u\) y \(-u\).
Para el coeficiente radial \(a_b\in\{4,1\}\),
\[
 \left\langle
 \frac{(\mathsf C^{\pm1}-I)a_b\rho}{3},a_b\rho
 \right\rangle=-a_b^2.
\]
La suma sobre los doce factores vuelve a dar \(-27\). La prueba del
\cref{km:estabilidad-vecino} se aplica componente a
componente.
\end{proof}

Para tipar la recalibración, sea
\(h\in\operatorname{Aut}(C_W)\) monomial: una permutación \(\sigma\) de
las doce coordenadas seguida de multiplicadores
\(\epsilon_b\in\{1,2\}\). Definimos su elevación reticular
\begin{equation}
 \widetilde h
 :=P_\sigma\circ
 \bigoplus_{b=0}^{11}\mathsf R^{\nu_b},
 \qquad
 \nu_b=
 \begin{cases}
 0,&\epsilon_b=1,\\
 1,&\epsilon_b=2.
 \end{cases}
 \label{paper:elevacion-monomial}
\end{equation}
Su acción en el discriminante es exactamente \(h\). Si el cambio transporta
simultáneamente código, bandera, origen incidencial, vector \(v\), vecino y
marcos, entonces
\begin{equation}
 g_{hu}=\widetilde h\,g_u\,\widetilde h^{-1}.
 \label{paper:naturalidad-isometrias}
\end{equation}
La matriz ternaria \(h\) no actúa sin elevación sobre \(A_2^{12}\) ni sobre
Leech.

